
% ============================================================
% Ejercicio 3
% ============================================================
\section{Ejercicio 3}

\begin{tcolorbox}[enunciado]
Describir e implemenar un algoritmo para generar certificados aleatoriios para ejemplares. El programa debe recibir como entrada (desde la línea de comandos):
\begin{itemize}
    \item nombre del archivo con la inforamción del ejemplar.
    \item nombre del archivo con el certificado a probrar.
\end{itemize}
El programa debe imprimir como salida:
\begin{itemize}
    \item Información general (similar a la que se imprimía en la práctica 1).
    \item Respuesta a la pregunta: ¿el ejemplar, con el certificado dado, satisface la condición de pertenencia al lenguaje correspondiente?
\end{itemize}
\end{tcolorbox}

\subsection{Algoritmo de Verificación}
El algoritmo primero revisa que la cardinalidad del certificado sea la misma con la cantidad del ejemplar, pues de otra manera el certificado no sería el correcto para el ejemplar.
Extrae W y V para tenerlos a la mano.
Ahora recorre de manera paralela Las tuplas del ejemplar junto con el certificado. Si el valor actual del certificado es 0, lo ignora y pasa al siguiente paso.
Si el valor actual del certificado es 1, entonces suma el peso y la ganancia actual, de igual manera verifica que el peso no exceda el de la mochila.
Si se llega a pasar, se detiene y devuelve no.
Si llega al final, verifica que haya pasado la meta, si no la paso, rechaza. Si el peso no se excede y cumple la meta, entonces devuelve si.

\begin{algorithm}[H]
\caption{Verificador de Certificados para el Problema de la Mochila}
\KwIn{Cadena binaria del ejemplar $I$, Cadena binaria del certificado $c$}
\KwOut{Respuesta sobre la pertenencia al lenguaje (SÍ / NO)}

\BlankLine
\tcp{Fase 1: Decodificación del Ejemplar}
$(P, W, V) \leftarrow \text{DesbinarizarEjemplar}(I)$ \tcp*{$P$: lista de tuplas $(w_i, v_i)$}
$n \leftarrow \text{Longitud}(P)$\;

\BlankLine
\tcp{Fase 2: Verificación de Estructura / Sintaxis}
\If{$\text{Longitud}(c) \neq n$}{
    \Return{\textbf{NO} (Error de longitud en el certificado)}\;
}

\BlankLine
\tcp{Fase 3: Recorrido en Paralelo y Evaluación Semántica}
$peso\_acumulado \leftarrow 0$\;
$ganancia\_acumulada \leftarrow 0$\;

\For{$i \leftarrow 1$ \KwTo $n$}{
    $bit \leftarrow c[i]$\;
    $(w_i, v_i) \leftarrow P[i]$\;
    
    \If{$bit = \text{'1'}{}$}{
        $peso\_acumulado \leftarrow peso\_acumulado + w_i$\;
        $ganancia\_acumulada \leftarrow ganancia\_acumulada + v_i$\;
        
        \tcp{Corte temprano en caso de exceso de peso}
        \If{$peso\_acumulado > W$}{
            \Return{\textbf{NO} (El peso acumulado excede la capacidad $W$)}\;
        }
    }
    \ElseIf{$bit \neq \text{'0'}{}$}{
        \Return{\textbf{NO} (Carácter inválido en el certificado)}\;
    }
}

\BlankLine
\tcp{Fase 4: Decisión Final de Pertenencia}
\If{$peso\_acumulado \le W$ \KwAnd $ganancia\_acumulada \ge V$}{
    \Return{\textbf{SÍ} (El certificado satisface las condiciones del lenguaje)}\;
}
\Else{
    \Return{\textbf{NO} (La ganancia acumulada no alcanza la meta $V$)}\;
}
\end{algorithm}